\chapter{Causal integral formulas}\label{app:causal-integrals}
For a unit impulse, $a=G_n$. For a force $f(t)$ of finite support,
$a=G_n*f$. Let $r(t)$ be the switched-sinusoid complex solution
\cref{eq:switched-solution}. A rectangular $N$-cycle burst is
\begin{equation}a_N^c(t)=r(t)-\Heaviside(t-NT)r(t-NT),\qquad T=2\pi/\Omega,
\end{equation}
because the phase delay $\ee^{-\ii\Omega NT}=1$. Direct evaluation at $t=0$
uses $C_++C_-=-H_n$ and
$s_+C_++s_-C_-=\ii\Omega H_n$, proving $a(0)=\dot a(0)=0$.

For an impulse spectrum $B(k)$ near a one-dimensional quadratic ZGV,
$\omega=\omega_0+D(k-k_0)^2/2$, the local integral is proportional to
\begin{equation}
B(k_0)\ee^{\ii(k_0x-\omega_0t)}
\int_{-\infty}^{\infty}\ee^{\ii[(k-k_0)x-Dt(k-k_0)^2/2]}\dd k,
\end{equation}
which is $O(t^{-1/2})$ for fixed $x$ if $D\ne0$ and $B(k_0)\ne0$. A two-dimensional
radial integral, a vanishing overlap, or a degenerate extremum changes the power.
